\documentclass[12pt,a4paper]{article}

\usepackage[brazil]{babel}
\usepackage{lmodern}
\usepackage{amsmath}
\usepackage{amsfonts}
\usepackage{listings}
% -----------------------------------------------------------------------
% PACOTES ESSENCIAIS
% -----------------------------------------------------------------------
\usepackage[utf8]{inputenc}        % Codificação do documento
\usepackage[T1]{fontenc}           % Codificação da fonte
\usepackage{microtype}             % Melhora a justificação
\usepackage{graphicx}              % Inclusão de figuras
\usepackage{indentfirst}           % Indenta o primeiro parágrafo
\usepackage{longtable}             % Tabelas longas
\usepackage{enumitem}              % Listas personalizáveis
\usepackage{hyperref}              % Hyperlinks
\usepackage{geometry}              % Ajuste de margens
\usepackage{titlesec}              % Customização dos títulos, se necessário
\usepackage[alf]{abntex2cite}

% -----------------------------------------------------------------------
% MARGENS (padrão ABNT é mais largas; aqui deixo como pediste)
% -----------------------------------------------------------------------
\geometry{
	left=2.5cm,
	right=2.5cm,
	top=2.5cm,
	bottom=2.5cm
}

% -----------------------------------------------------------------------
% CONFIGURAÇÕES DO HYPERREF
% -----------------------------------------------------------------------
\hypersetup{
	colorlinks=true,
	linkcolor=black,
	citecolor=black,
	urlcolor=blue
}

% -----------------------------------------------------------------------
% TÍTULO, AUTOR E DATA
% -----------------------------------------------------------------------
\title{
	Ontological Clarity Index (OCI) e Ontological Compatibility Score (OCS):\\
	Métricas para Convergência e Integração de Sistemas Complexos
}

\author{José Pedro Trindade\thanks{jose.pereira-trindade@embrapa.br}}
\date{\today}

\begin{document}

	\maketitle
	\tableofcontents



	\section{Introdução}

	Sistemas complexos contemporâneos, especialmente aqueles compostos por múltiplos projetos, módulos ou iniciativas independentes, apresentam um desafio recorrente: a dificuldade de integração, convergência e evolução coordenada.

	Na ausência de critérios estruturados, decisões de integração são frequentemente baseadas em similaridade superficial, proximidade histórica ou julgamento subjetivo. Isso resulta em redundância, fragmentação e aumento do custo de manutenção.

	Este trabalho propõe uma abordagem baseada em explicitação ontológica, introduzindo duas métricas: o \textit{Ontological Clarity Index (OCI)} e o \textit{Ontological Compatibility Score (OCS)}, com o objetivo de tornar sistemas comparáveis, integráveis e analisáveis de forma sistemática.

Diferentemente de abordagens tradicionais que focam em representação ou alinhamento, esta proposta enfatiza a mensuração da qualidade ontológica e da compatibilidade estrutural entre sistemas, estabelecendo uma ponte entre modelagem conceitual e avaliação quantitativa.

\section{Trabalhos Relacionados}

A avaliação e alinhamento de ontologias é um tema amplamente estudado nas áreas de engenharia de conhecimento e web semântica.

Gruber (1993) define ontologias como especificações explícitas de conceitualizações, estabelecendo a base para sua formalização computacional.

Guarino (1998) amplia essa definição ao introduzir uma distinção entre níveis ontológicos, enfatizando o papel da ontologia formal na modelagem conceitual rigorosa.

No contexto da web semântica, ontologias são frequentemente representadas em linguagens formais como OWL (Web Ontology Language), permitindo inferência automática e interoperabilidade entre sistemas (W3C, 2004).

A literatura sobre alinhamento de ontologias (ontology matching) aborda a identificação de correspondências entre entidades de diferentes ontologias (Euzenat e Shvaiko, 2013), sendo uma área diretamente relacionada ao conceito proposto de compatibilidade ontológica.

Entretanto, a maioria dos trabalhos existentes foca em algoritmos de correspondência ou integração, enquanto há menor ênfase na definição de métricas estruturais que permitam avaliar, de forma sistemática, o grau de clareza interna e compatibilidade entre ontologias.

Neste contexto, o presente trabalho propõe os índices OCI e OCS como métricas complementares, visando preencher essa lacuna ao oferecer uma abordagem quantitativa para análise ontológica.

\section{Fundamentação Ontológica}

O LabGestao configura-se como um ambiente laboratorial de análise, caracterização e avaliação de sistemas. Essa caracterização não se refere a um laboratório experimental clássico, mas a um contexto analítico estruturado no qual práticas, processos, ferramentas e métricas são articulados para observar, descrever, comparar e avaliar outros sistemas.

Nesse sentido, o LabGestao não opera apenas como um sistema de gestão de projetos. Ele funciona também como um dispositivo analítico: um ambiente no qual princípios ontológicos e estruturais são manifestos, operacionalizados e aplicados sobre repositórios observados. O conjunto formado por inventário, governança, métricas, relatórios, recomendações e perfis ontológicos permite que outros sistemas sejam lidos segundo critérios explícitos de clareza, compatibilidade, maturidade e prontidão de integração.

No caso específico de OCI e OCS, as métricas não devem ser interpretadas como indicadores isolados. Elas expressam formas formais de observação: o OCI explicita como um sistema define entidades, relações, validade e identidade; o OCS explicita em que medida dois sistemas podem ser comparados, alinhados ou integrados a partir desses mesmos critérios. Assim, o LabGestao pode ser compreendido como um laboratório virtual de análise ontológica e estrutural de sistemas, pois incorpora em seu próprio modo de operação os princípios que utiliza para avaliar outros sistemas.

No contexto de sistemas computacionais, uma ontologia é definida como uma estrutura formal:

\[
\mathcal{O} = (E, R, V, I)
\]

onde:

\begin{itemize}
	\item \(E\) é um conjunto finito de entidades;
	\item \(R \subseteq E \times E\) é um conjunto de relações tipadas entre entidades;
	\item \(V\) é um conjunto de regras de validade (predicados sobre \(E\) e \(R\));
	\item \(I\) é um conjunto de regras de identidade (funções ou critérios que distinguem elementos de \(E\)).
\end{itemize}

Essa estrutura permite tratar sistemas como objetos comparáveis, superando descrições informais.

Uma ontologia mínima deve responder a quatro questões fundamentais:

\begin{itemize}
	\item O que existe no sistema? (Entidades)
	\item Como essas coisas se relacionam? (Relações)
	\item O que pode ser considerado válido? (Regras de validade)
	\item Como distinguir e identificar elementos? (Identidade)
\end{itemize}

No contexto do LabGestao, essa ontologia mínima convive com artefatos de engenharia que não são ontologia em si, mas funcionam como fontes de explicitação, evidência e rastreabilidade:

\begin{itemize}
	\item \textbf{DDD}: principal fonte para identificar entidades, relações, fronteiras e linguagem ubíqua;
	\item \textbf{ADR}: explicita decisões arquiteturais e restrições relevantes para a interpretação ontológica;
	\item \textbf{DAI}: registra decisões, ações e impedimentos que contextualizam a evolução estrutural do sistema;
	\item \textbf{CI}: não define ontologia diretamente, mas fornece evidência operacional de validação e disciplina de execução.
\end{itemize}

\subsection{Definição Operacional de Entidade}

A definição de entidade é central para a consistência do \textit{Ontological Clarity Index} (OCI) e do \textit{Ontological Compatibility Score} (OCS). Embora a ontologia proposta represente \(E\) como um conjunto finito de entidades, essa formulação, por si só, ainda é insuficiente para garantir comparabilidade, reprodutibilidade e rigor analítico. Sem uma definição operacional mais precisa, diferentes avaliadores podem classificar como entidade elementos heterogêneos, produzindo variações artificiais no cálculo das métricas.

Em termos práticos, quase qualquer componente de um sistema poderia ser tratado como entidade: objetos do domínio, módulos técnicos, documentos, eventos, atores, métricas, processos ou artefatos de governança. Essa elasticidade conceitual compromete a robustez da avaliação, pois permite inflacionar ou reduzir o conjunto \(E\) sem critério ontológico suficientemente explícito.

Diante disso, propõe-se a seguinte definição operacional:

\begin{quote}
	Uma entidade é um elemento ontologicamente relevante do domínio que possui identidade distinguível, participa de relações estruturais, pode ser submetido a regras de validade e apresenta papel persistente na interpretação do sistema.
\end{quote}

Essa definição permite distinguir entidades propriamente ditas de elementos auxiliares, contextuais ou instrumentais. Em termos formais, pode-se representar o conjunto de entidades como:

\[
E = \{x \mid x \text{ possui identidade, papel estrutural e relevância ontológica no domínio}\}
\]

\subsubsection{Critérios para Reconhecimento de Entidades}

Para fins de análise, um elemento só deve ser incorporado ao conjunto \(E\) quando satisfizer, ao menos de modo predominante, os seguintes critérios:

\begin{itemize}
	\item \textbf{Identidade própria}: o elemento pode ser distinguido de outros elementos do mesmo tipo;
	\item \textbf{Persistência conceitual}: o elemento mantém significado estável ao longo do sistema, não sendo apenas uma ocorrência ocasional;
	\item \textbf{Capacidade relacional}: o elemento participa de relações estruturalmente relevantes;
	\item \textbf{Validade avaliável}: o elemento pode ser submetido a regras, restrições ou critérios de consistência;
	\item \textbf{Relevância de domínio}: o elemento é constitutivo do domínio analisado, e não apenas suporte técnico ou documental.
\end{itemize}

Esses critérios reduzem arbitrariedades e tornam mais transparente a distinção entre o que pertence ao núcleo ontológico do sistema e o que atua apenas como evidência, contexto ou mecanismo de suporte.

\subsubsection{Camadas de Classificação Ontológica}

Para reforçar a precisão analítica, recomenda-se distinguir pelo menos quatro camadas:

\begin{enumerate}
	\item \textbf{Entidades de domínio}: representam aquilo que o sistema reconhece como existente no problema tratado, como \textit{Project}, \textit{Metric}, \textit{Event}, \textit{Actor} ou \textit{Recommendation};

	\item \textbf{Entidades técnicas}: elementos existentes na infraestrutura ou implementação, como \textit{Repository}, \textit{API}, \textit{Database}, \textit{Pipeline} ou \textit{CIJob};

	\item \textbf{Artefatos de governança e documentação}: elementos como DDD, ADR, DAI, políticas e esquemas, que funcionam como fontes de explicitação e evidência, mas que não devem ser considerados automaticamente entidades ontológicas do domínio;

	\item \textbf{Eventos ou estados derivados}: elementos como \textit{ValidationFailure}, \textit{Deployment} ou \textit{IntegrationDecision}, cuja classificação como entidade depende de seu papel estrutural e persistência conceitual.
\end{enumerate}

Essa diferenciação é especialmente importante porque diversos artefatos de engenharia contribuem para a explicitação da ontologia sem, contudo, constituírem ontologia em si. DDD, ADR, DAI e CI, por exemplo, funcionam como fontes de evidência, rastreabilidade e confiança, mas não devem ser incorporados ao conjunto \(E\) sem justificação conceitual explícita.

\subsubsection{Implicações para o OCI e o OCS}

A explicitação mais rigorosa do conceito de entidade possui impacto direto nas métricas propostas.

Primeiro, melhora a \textbf{comparabilidade}, evitando que projetos distintos sejam avaliados a partir de critérios incompatíveis sobre o que conta como entidade.

Segundo, aumenta a \textbf{reprodutibilidade}, pois reduz a dependência de julgamentos implícitos na construção de \(E_{expected}\) e \(E_{detected}\).

Terceiro, fortalece a \textbf{consistência do OCS}, uma vez que o alinhamento entre ontologias depende, em primeiro lugar, da qualidade da delimitação de suas entidades. Se os conjuntos de entidades forem formados por critérios heterogêneos, o valor de \(M_E\) poderá refletir apenas correspondências superficiais, e não compatibilidade ontológica efetiva.

Desse modo, a clareza na definição de entidade não deve ser tratada como detalhe metodológico secundário, mas como fundamento da validade analítica do framework proposto.

\subsubsection{Nota Metodológica}

Nem todo componente mencionado em artefatos de engenharia constitui uma entidade ontológica. Documentos, decisões, mecanismos técnicos e evidências operacionais devem ser distinguidos de entidades de domínio, ainda que possam servir como fontes relevantes para sua inferência. Essa distinção é necessária para evitar que o OCI seja reduzido a uma contagem heurística de termos recorrentes, preservando seu caráter de métrica voltada à explicitação ontológica do sistema.

\section{Ontological Clarity Index (OCI)}

\subsection{Definição}

O OCI é definido como uma função:

\[
OCI : \mathcal{O} \rightarrow [0,1]
\]

que mede o grau de explicitação ontológica de um sistema.

\subsection{Componentes}

Dada uma ontologia \(\mathcal{O} = (E, R, V, I)\), definimos quatro funções de avaliação:

\begin{itemize}
	\item \(C_E(E)\): cobertura e definição das entidades;
	\item \(C_R(R)\): explicitação e completude das relações;
	\item \(C_V(V)\): formalização das regras de validade;
	\item \(C_I(I)\): clareza dos critérios de identidade.
\end{itemize}

Cada componente é normalizado no intervalo \([0,1]\), de acordo com critérios explícitos de avaliação (por exemplo: completude, ausência de ambiguidade e verificabilidade).

\subsection{Formulação Operacional}

Para garantir reprodutibilidade, os componentes são definidos como:

\[
C_E = \frac{|E_{detected} \cap E_{expected}|}{|E_{expected}|}
\]

\[
C_R = \frac{|R_{detected} \cap R_{expected}|}{|R_{expected}|}
\]

\[
C_V \in \{0, 0.5, 1\}
\quad
C_I \in \{0, 0.5, 1\}
\]

\subsection{Definição de Referência}

Os conjuntos \(E_{expected}\) e \(R_{expected}\) são definidos por:

\begin{itemize}
	\item ontologia canônica do domínio;
	\item artefatos de modelagem do projeto, especialmente DDD;
	\item baseline manual do projeto.
\end{itemize}

Em termos operacionais, o DDD tende a ser a principal fonte para definição de \(E_{expected}\) e \(R_{expected}\), enquanto ADR e DAI atuam como artefatos auxiliares de contextualização e rastreabilidade.

\subsection{Observação}

A normalização exige definição prévia de um espaço de referência (por exemplo, domínio esperado ou requisitos do sistema), evitando arbitrariedade na avaliação.

\section{Ontological Compatibility Score (OCS)}

\subsection{Definição}

O OCS é definido como uma função:

\[
OCS : \mathcal{O}_A \times \mathcal{O}_B \rightarrow [0,1]
\]

que mede o grau de compatibilidade entre duas ontologias.

\subsection{Componentes}

Sejam duas ontologias:

\[
\mathcal{O}_A = (E_A, R_A, V_A, I_A), \quad
\mathcal{O}_B = (E_B, R_B, V_B, I_B)
\]

Definimos:

\begin{itemize}
	\item \(M_E(E_A, E_B)\): grau de alinhamento entre entidades;
	\item \(M_R(R_A, R_B)\): compatibilidade estrutural das relações;
	\item \(M_V(V_A, V_B)\): consistência entre regras de validade;
	\item \(M_I(I_A, I_B)\): compatibilidade dos critérios de identidade.
\end{itemize}

\subsection{Formulação}
Para garantir consistência com os axiomas:

\[
OCS =
\begin{cases}
	0, & \text{se } M_E = 0 \\
	\frac{M_E + M_R + M_V + M_I}{4}, & \text{caso contrário}
\end{cases}
\]

\subsection{Observação}

Essa definição garante que não há compatibilidade ontológica sem interseção semântica mínima entre entidades, evitando falsos positivos.

\subsection{Interpretação}

\begin{itemize}
	\item Alto (>$0.7$): forte potencial de integração
	\item Médio ($0.4$--$0.7$): requer alinhamento estrutural
	\item Baixo (<$0.4$): baixa compatibilidade ontológica
\end{itemize}

A função \(M_E\) pode ser interpretada como um caso particular de alinhamento ontológico, enquanto \(M_R\), \(M_V\) e \(M_I\) ampliam essa noção ao considerar não apenas correspondência estrutural, mas também consistência semântica e validade lógica.

Dessa forma, o OCS generaliza abordagens clássicas de ontology matching ao incorporar múltiplas dimensões de compatibilidade.

\section{Propriedades e Axiomas}

Para garantir consistência e evitar arbitrariedade, propõem-se os seguintes axiomas:

\begin{itemize}
	\item A1: Ambiguidade reduz \(OCI\)
	\item A2: Ausência de identidade limita \(OCI\)
	\item A3: \(OCS(\mathcal{O}, \mathcal{O}) = 1\)
	\item A4: Se \(M_E = 0\), então \(OCS = 0\)
	\item A5: Conflitos reduzem \(OCS\)
\end{itemize}

\section{Fator de Confiança}

Define-se:

\[
confidence \in [0,1]
\]

baseado em:

\begin{itemize}
	\item origem dos dados (código, documentação, heurística)
	\item presença de artefatos formais (DDD, ADR, contratos, schemas)
	\item consistência interna detectada
	\item evidência operacional de validação (incluindo CI)
\end{itemize}

O fator de confiança não altera a definição ontológica do sistema. Seu papel é ajustar a confiança na medição realizada. Assim, DDD, ADR e DAI funcionam como fontes de estrutura e rastreabilidade, enquanto CI funciona principalmente como evidência de validação operacional.

\[
OCI_{adj} = OCI \times confidence
\]

\[
OCS_{adj} = OCS \times confidence
\]

\section{Operacionalização Heurística (LabGestao)}

A implementação prática utiliza:

\begin{itemize}
	\item extração textual (README, DDD, ADR, documentação arquitetural e código fonte)
	\item identificação de entidades via padrões semânticos
	\item detecção de relações via verbos estruturais
	\item validade por presença de regras explícitas
	\item identidade por identificadores únicos
\end{itemize}

\subsection{Papel dos Artefatos de Engenharia}

Na operacionalização atual do LabGestao:

\begin{itemize}
	\item \textbf{DDD} tem peso prioritário para inferir entidades, relações e fronteiras esperadas;
	\item \textbf{ADR} reforça justificativas estruturais, restrições e decisões de modelagem;
	\item \textbf{DAI} não define ontologia diretamente, mas ajuda a interpretar estágio e contexto de evolução do projeto;
	\item \textbf{CI} não aumenta diretamente \(OCI\), mas pode elevar o fator de confiança ao demonstrar validação reprodutível.
\end{itemize}

\subsection{Limitação}

Essa abordagem pode superestimar entidades em projetos com alta densidade textual.

Além disso, a heurística atual depende de artefatos efetivamente versionados. Repositórios com DDD, ADR ou DAI ausentes no código-fonte, mas presentes apenas em conhecimento tácito ou documentos externos, tendem a receber escores subestimados.


\section{Recommendation Engine Ontológico}

Define-se:

\[
integration\_readiness = OCI \times OCS \times governance
\]

onde \(governance \in [0,1]\) representa um fator derivado de evidências estruturais e operacionais, incluindo DDD, ADR, DAI, políticas, contratos e sinais de validação.

\subsection{Regras de Decisão}

\begin{itemize}
	\item OCI baixo → priorizar estrutura ontológica
	\item OCS alto + OCI baixo → evitar integração prematura
	\item OCS alto + OCI alto → integrar
	\item OCS médio → alinhar
	\item OCS baixo → manter separado
\end{itemize}

Esse motor não deve ser interpretado como regra mecânica suficiente. Sua função é orientar movimentos de evolução ontológica do sistema, combinando clareza interna, compatibilidade externa e contexto operacional.


\section{Exemplo de Cálculo}

Considere uma ontologia com:

\begin{itemize}
	\item 3 entidades definidas de um total esperado de 4 → \(C_E = 0.75\)
	\item 2 relações explicitadas de 3 esperadas → \(C_R = 0.67\)
	\item regras de validade parcialmente definidas → \(C_V = 0.5\)
	\item identidade bem definida → \(C_I = 1.0\)
\end{itemize}

Então:

\[
OCI = \frac{0.75 + 0.67 + 0.5 + 1.0}{4} = 0.73
\]

\subsection{Exemplo de Compatibilidade}

A compatibilidade entre ontologias pode ser estimada via:

\begin{itemize}
	\item interseção normalizada de entidades;
	\item correspondência de relações;
	\item consistência entre regras de validade e identidade.
\end{itemize}


	\section{Aplicação em Sistemas de Engenharia}

	As métricas propostas podem ser aplicadas em:

	\begin{itemize}
		\item Portfólios de projetos
		\item Sistemas modulares
		\item Arquiteturas distribuídas
	\end{itemize}

	Quando integradas a estruturas de análise como grafos, permitem identificar:

	\begin{itemize}
		\item Redundâncias
		\item Oportunidades de integração
		\item Gargalos estruturais
	\end{itemize}

	\section{Exemplo Aplicado}

	Considere dois sistemas:

	\begin{itemize}
		\item Sistema A: foco em gestão de projetos
		\item Sistema B: foco em avaliação e métricas
	\end{itemize}

	Se ambos compartilham entidades como \textit{Project} e \textit{Metric}, mas diferem em regras de validade, o OCS indicará compatibilidade parcial, sugerindo necessidade de alinhamento antes da integração.

\section{Resultados Experimentais}

A análise de 75 repositórios apresentou:

\begin{itemize}
	\item Média OCI: 42.6
	\item Alta frequência de:
	\begin{itemize}
		\item entidades bem definidas
		\item ausência de relações
		\item ausência de identidade
	\end{itemize}
\end{itemize}

\subsection{Interpretação}

Os resultados indicam um padrão de:

\begin{quote}
	estruturação ontológica parcial
\end{quote}

\subsection{Limitação}

Os valores dependem do método heurístico e devem ser interpretados como aproximação.

	\section{Discussão}
As métricas propostas avaliam:

\begin{quote}
	grau de explicitação ontológica e não qualidade funcional do sistema
\end{quote}

Principais limitações:

\begin{itemize}
	\item dependência de baseline ontológico
	\item sensibilidade à qualidade da documentação
	\item natureza heurística da detecção
	\item necessidade de calibrar separadamente ontologia, governança e evidência operacional
\end{itemize}

Apesar disso, os resultados demonstram capacidade de:

\begin{itemize}
	\item identificar fragilidade estrutural
	\item sugerir decisões de integração
	\item revelar padrões sistêmicos
\end{itemize}

Em particular, a combinação entre DDD, ADR, DAI e CI não deve ser interpretada como substituto da ontologia. O valor desses artefatos está em aumentar a capacidade de explicitação, comparação e confiança da análise. O DDD tende a influenciar mais diretamente o conteúdo ontológico; ADR e DAI reforçam contexto e rastreabilidade; CI reforça evidência operacional.
		\section{Conclusão}

	Este trabalho propõe um framework ontológico para avaliação de sistemas complexos, baseado nos índices OCI e OCS.

	Sua principal contribuição é a transformação de estruturas conceituais em elementos mensuráveis, permitindo decisões mais informadas sobre integração, evolução e descarte de sistemas.

	Ao tornar ontologias comparáveis, abre-se caminho para arquiteturas mais coerentes, adaptativas e integráveis.

\begin{thebibliography}{9}

	\bibitem{gruber1993}
	GRUBER, T. R. A translation approach to portable ontology specifications.
	\textit{Knowledge Acquisition}, 1993.

	\bibitem{guarino1998}
	GUARINO, N. Formal ontology in information systems.
	In: \textit{Proceedings of FOIS}, 1998.

	\bibitem{euzenat2013}
	EUZENAT, J.; SHVAIKO, P.
	\textit{Ontology Matching}.
	Springer, 2013.

	\bibitem{w3cowl}
	W3C. OWL Web Ontology Language Reference.
	2004.

	\bibitem{bernerslee2001}
	BERNERS-LEE, T.; HENDLER, J.; LASSILA, O.
	The Semantic Web.
	\textit{Scientific American}, 2001.

\end{thebibliography}

\end{document}
